% Folder 70, batch 4 — archivists' pages 61–80.
% Pass: Opus 5.5 (claude-opus-5-5), 2026-09-29 — first pass, unchecked against the pages by a human.
%
% Skipped as carrying no mathematics: 63 (a blank verso showing the ink of
% p. 62 through the paper), 72 (blank coloured separator sheet), 77 and 79
% (typed administrative versos). Page 73 is a cover inscribed in his hand; its
% words become the heading before p. 74 and it gets no page marker.
% Notation fixed for the batch: \sigma_0,\sigma_1,\sigma_2 the three reflections
% (generators), \rho_s, \rho_f their products, s_0..s_3, s'_2, s'_3, s''_3 the
% vertices, e_i = s_i - s_0, \alpha' the conjugate of \alpha (\alpha+\alpha'=-1),
% the antipodism he writes as an underlined a is set \mathbf{a}.
\documentclass[11pt,a4paper]{article}
\input{../preamble/grothendieck}

\folder{70}
\batch{4}
\pages{61}{80}
\foldertitle{Réalisations géométriques de structures combinatoires (n-polyèdres, n-hyperpolyèdres [2-polyèdres réguliers]…) (Vieilles rédactions) : notes manuscrites (s.d.).}
\dating{[à partir de 1976]}
\watermark{Édition de démonstration}

\begin{document}

\section{Le groupe cartographique et $\mathrm{Sl}(2,\mathbb{Z})$}

\page{61}
\note{la page s'ouvre au milieu d'un calcul : $\sigma$, $\rho_s$, $\rho_f$ et $\Gamma^0_{\infty,3}$ sont déjà en usage, l'argument vient donc des pages qui précèdent ce lot.}

\[
1 \to R_n \to \mathrm{Sl}(2,\mathbb{Z})/\pm 1 \to \mathrm{Sl}(2,\mathbb{Z}_2) \to 0
\]
\note{l'indice de $R$ est douteux (\uncertain{$n$} ou $2$) ; le « $/\pm 1$ » est écrit par-dessus une rature.}
avec $R_n = R(\infty,3;2,3)$ engendré par $\rho_s^2 = (\sigma\rho_f^{-1})^2$ ; $\mathrm{Sl}(2,\mathbb{Z})/\pm 1 \simeq \Gamma^0_{\infty,3}$, $\mathrm{Sl}(2,\mathbb{Z}_2) \simeq \mathfrak{S}_3 \simeq \Gamma_{2,3}$.

\[
\begin{cases}
\rho_s = \sigma_0\sigma_1 \\
\rho_f = \sigma_1\sigma_2 \\
\rho_s\rho_f = \sigma_0\sigma_2 = \sigma
\end{cases}
\qquad \rho_s = \sigma\rho_f^{-1},
\qquad \boxed{\rho_f^3 = 1,\ \ \sigma^2 = 1}
\]

\[
\frac{az+b}{cz+d} = z, \qquad cz^2 + (d-a)z - b = 0, \qquad z = x+iy \ (y \neq 0)
\]
\[
\begin{cases}
c(x^2-y^2) + (d-a)x - b = 0 \\
\bigl[2cx + (d-a)\bigr]y = 0 \iff 2cx + (d-a) = 0
\end{cases}
\Longleftrightarrow
\begin{cases}
d-a = -2cx \\
c(x^2+y^2) + b = 0
\end{cases}
\]
\[
\begin{cases}
d = a - c\,\mathrm{Tr}(z) \\
b = -c\,N(z)
\end{cases}
\]
\[
\begin{pmatrix} a & b \\ c & d \end{pmatrix}
= \begin{pmatrix} a & -cN(z) \\ c & a - c\,\mathrm{Tr}\,z \end{pmatrix}
\]
\[
\det = a^2 - ac\,\mathrm{Tr}\,z + c^2N(z) = (a-cz)(a-c\bar z) = N(a-cz) = 1
\]

1°) $z = i$, $\mathrm{Tr}\,z = 0$, $N(z) = 1$ ($z$ racine primitive 4\supplied{\textsuperscript{e}} de 1)
\[
G_z = G_i = \left\{ \begin{pmatrix} a & -c \\ c & a \end{pmatrix} \,\middle|\, a^2+c^2 = 1 \right\}
\]
élément d'ordre 4 :
\[
\boxed{\sigma = \begin{pmatrix} 0 & -1 \\ 1 & 0 \end{pmatrix}}
\qquad \sigma^4 = \mathrm{id}, \qquad \boxed{\sigma^2 = -1}
\]

2°) $z = \frac12 + i\frac{\sqrt3}{2} = \zeta$ ($z$ racine primitive 6\supplied{\textsuperscript{e}} de 1), $\mathrm{Tr}\,z = 1$, $Nz = 1$
\[
G_z = G_\zeta = \left\{ \begin{pmatrix} a & -c \\ c & a-c \end{pmatrix} \,\middle|\, a^2 - ac + c^2 = 1 \right\},
\qquad N(a - c\zeta) = 1
\]
\[
\boxed{\rho_f = \begin{pmatrix} 0 & 1 \\ -1 & 1 \end{pmatrix}}
\qquad \rho^6 = \mathrm{id}, \qquad \boxed{\rho^3 = -1}
\]
\note{les signes des deux coefficients de droite de $\rho_f$ sont repassés et en partie noyés d'encre ; un symbole biffé précède $\rho_f$.}
\[
\boxed{\rho'_s = \sigma\rho_f = \sigma_2\rho_f\sigma_2^{-1} = \begin{pmatrix} 1 & -1 \\ 0 & 1 \end{pmatrix}}
\qquad \rho'^{\,n}_s = \begin{pmatrix} 1 & -n \\ 0 & 1 \end{pmatrix}
\]

\page{62}
\[
\Gamma^0_{\infty,3} \simeq \mathrm{Sl}(2,\mathbb{Z})/\pm 1
\]
\[
\Gamma^0_{n,3} \xrightarrow{\ \text{épi}\ } \mathrm{Sl}(2,\mathbb{Z}_n)/\pm 1 \quad (n \geq 2) \quad (\text{iso \emph{ssi} } n \leq 5)
\]
\marginal{i.e. ssi $\Gamma^0_n$ fini}
\struck{$\simeq \mathrm{Gl}(2,\mathbb{Z}_n)$}

l'étudier pour \emph{$n = 2, 3, 4, 5, 6$}

\[
\begin{array}{c|c|c|l}
n & \text{ordre } \Gamma^0_{n,3} & \text{ordre } \mathrm{Sl}(2,\mathbb{Z}_n)/\pm 1 & \\ \hline
2 & (\mathfrak{S}_3)\ 6 & 6 & \text{cas 3-diédral} \\
3 & (\mathfrak{A}_4)\ 12 & 12 & \text{tétraèdre} \\
4 & (\mathfrak{S}_4)\ 24 & 24 & \text{cube orienté} \\
5 & (\mathfrak{A}_5)\ 60 & 60 & \text{icos. orienté} \\
6 & (\mathbb{Z}_6\cdot\mathbb{Z}^2)\ \infty & 72 & \text{pavage par triangles orienté du plan}
\end{array}
\]
\note{la dernière case (« \uncertain{pavage par triangles orienté du plan} ») est très serrée ; « cas 3-diédral » est aussi une lecture probable plutôt que sûre. Un mot biffé \struck{\ill{}} précède le tableau.}

\[
\mathrm{card}\,\mathrm{Sl}(2,\mathbb{F}_q) = (q+1)q(q-1) = q(q^2-1)
\]
\[
\mathrm{card}\,\mathrm{Sl}(2,\mathbb{Z}/p^\nu\mathbb{Z}) = \bigl(\mathrm{card}\,\mathrm{Sl}(2,\mathbb{F}_p)\bigr)\,p^{3(\nu-1)}
\]

\section{Les trois réflexions en coordonnées}

\page{64}
$\sigma_0\ \sigma_1\ \sigma_2$, \quad $r = (s_0, a_0, f_0)$

\[
\begin{array}{lll}
\sigma_0 s_0 = s_1 & \sigma_0 a_0 = a_0 & \sigma_0 f_0 = f_0 \\
\sigma_1 s_0 = s_0 & \sigma_1 a_0 = \{s_0, s_2\} & \sigma_1 f_0 = f_0 \\
\sigma_2 s_0 = s_0 & \sigma_2 a_0 = a_0 & \sigma_2 f_0 = \{s_0, s_1, s_3, \ldots\}
\end{array}
\qquad
\begin{array}{l}
a_0 = \{s_0, s_1\} \\
f_0 = \{s_0, s_1, s_2, \ldots\}
\end{array}
\]

\[
\sigma_0 s_0 = s_1, \qquad \sigma_1 s_1 = s_2, \qquad \sigma_2 s_2 = s_3
\]

\drawing{Figure autour du drapeau $(s_0, a_0)$ : l'arête $s_0 s_1$ (marquée $a_0$), le triangle $s_0 s_1 s_2$ et ses images ; les sommets $s_2$, $s_3$, $s'_2$, $s'_3$ et, en pointillé à gauche, $s''_3$ ; les réflexions $\sigma_0$, $\sigma_1$, $\sigma_2$ sont marquées sur les droites qu'elles fixent.}

\[
\begin{cases}
e_1 = s_1 - s_0 \\
e_2 = s_2 - s_0 \\
e_3 = s_3 - s_0
\end{cases}
\]

\[
\sigma_1 \begin{cases} s_0 \mapsto s_0 \\ s_1 \mapsto s_2 \\ s_2 \mapsto s_1 \\ s_3 \mapsto s''_3 \end{cases}
\qquad
\sigma_2 \begin{cases} s_0 \mapsto s_0 \\ s_1 \mapsto s_1 \\ s_2 \mapsto s_3 \\ s_3 \mapsto s_2 \end{cases}
\qquad
\sigma_0 \begin{cases} s_0 \mapsto s_1 \\ s_1 \mapsto s_0 \\ s_2 \mapsto s'_2 \\ s_3 \mapsto s'_3 \end{cases}
\]

\[
\boxed{\mathrm{Pl}(s_0, s_1, s_2) = \mathrm{Pl}(s_0, s_1, s'_2)}
\qquad
\boxed{\mathrm{Pl}(s_0, s_1, s_3) = \mathrm{Pl}(s_0, s_1, s'_3)}
\]
\note{dans le premier cadre, une première écriture \struck{$(s_0, s_1, s_2)$} est biffée.}

Tableau de $\sigma_0\sigma_2 = \sigma_2\sigma_0$ (encadré) :
\[
\begin{array}{l|ll}
s_0 \mapsto & s_1 & s_1 \\
s_1 \mapsto & s_0 & s_0 \\
s_2 \mapsto & s'_3 & \sigma_2 s'_2 \\
s_3 \mapsto & s'_2 & \sigma_2 s'_3
\end{array}
\]
\[
\boxed{\sigma_0^2 = 1} \quad \boxed{\sigma_1^2 = 1} \quad \sigma_2^2 = 1 \text{ automatique}
\]
\[
\boxed{\begin{array}{l} \sigma_0 s'_2 = s_2 \\ \sigma_0 s'_3 = s_3 \end{array}}
\qquad
\boxed{\sigma_1 s''_3 = s_3}
\qquad
\boxed{\sigma_2 s'_2 = s'_3}, \quad \sigma_2 s'_3 = s'_2
\]
\note{des flèches doubles relient chacune de ces relations encadrées à la relation d'involution ou de commutation qu'elle exprime.}

\[
s'_2 = s_0 + \lambda e_1 + \mu e_2 = (\lambda, \mu, 0),
\qquad
s'_3 = s_0 + \lambda e_1 + \mu e_3 = (\lambda, 0, \mu)
\]
\[
\Bigl[\ \underbrace{\mu = 1}_{?},\ \lambda = 1 + \beta \ldots\ \Bigr]
\]
\[
s''_3 = s_3 + \xi e_1 + \eta e_2 + \zeta e_3 = s_0 + \xi e_1 + \eta e_2 + (1+\zeta)e_3 = (\xi, \eta, 1-\zeta)
\]
\[
\Bigl[\ \underbrace{\xi = -\eta}_{?},\ \eta = 1 + \alpha,\ \underbrace{\zeta = 0}_{?}\ \Bigr]
\]
\note{« $1-\zeta$ » dans la dernière coordonnée, contre « $1+\zeta$ » dans le membre du milieu : tel quel sur la page.}

\drawing{Petit schéma : le carré $s_0 s_1 s'_2 s_2$ vu de face, avec $s_3$, $s'_3$ en dessous et $s''_3$ à gauche relié en pointillé.}

\[
\begin{cases}
\sigma_0 s_0 = s_1 = (1,0,0) \\
\sigma_0 e_1 = -e_1 \\
\sigma_0 e_2 = (\lambda-1)e_1 + \mu e_2 \\
\sigma_0 e_3 = (\lambda-1)e_1 + \mu e_3
\end{cases}
\qquad
\sigma_0 = \begin{pmatrix}
-1 & \lambda-1 & \lambda-1 & 1 \\
0 & \mu & 0 & 0 \\
0 & 0 & \mu & 0 \\
0 & 0 & 0 & 1
\end{pmatrix}
\]

\page{65}
\[
\boxed{\sigma_0(x,y,z) = \bigl(1 - x + (\lambda-1)y + (\lambda-1)z,\ \mu y,\ \mu z\bigr)}
\]
\[
\sigma_0(s'_2) = \sigma_0(\lambda,\mu,0) = \bigl(1-\lambda+(\lambda-1)\mu,\ \mu^2,\ 0\bigr) \overset{?}{=} s_2 = (0,1,0)
\]
i.e.
\[
\boxed{\mu^2 = 1 \qquad (\lambda-1)(\mu-1) = 0}
\]
$\sigma_0(s'_3) = s_3$ i.e. $\sigma_2\sigma_0(s'_3) = \sigma_2 s_3$, où $\sigma_2\sigma_0 = \sigma_0\sigma_2$, $\sigma_0\sigma_2 s'_3 = \sigma_0 s'_2$ et $\sigma_2 s_3 = s_2$.

\[
\begin{cases}
\sigma_1 s_0 = s_0 = (0,0,0) \\
\sigma_1 e_1 = e_2 \\
\sigma_1 e_2 = e_1 \\
\sigma_1 e_3 = \xi e_1 + \eta e_2 + (1+\zeta)e_3
\end{cases}
\qquad
\sigma_1 = \begin{pmatrix}
0 & 1 & \xi & 0 \\
1 & 0 & \eta & 0 \\
0 & 0 & 1+\zeta & 0 \\
0 & 0 & 0 & 1
\end{pmatrix}
\]
\note{un mot biffé \struck{\ill{}} précède l'accolade.}
\[
\boxed{\sigma_1(x,y,z) = \bigl(y + \xi z,\ x + \eta z,\ (1+\zeta)z\bigr)}
\]
\[
\sigma_1(s''_3) = \sigma_1(\xi, \eta, 1-\zeta) = \Bigl(\eta + \xi(1-\zeta),\ \xi + \eta(1-\zeta),\ \underbrace{(1+\zeta)(1-\zeta)}_{1-\zeta^2}\Bigr)
\overset{?}{=} s_3 = (0,0,1)
\]
\[
\zeta^2 = 0, \qquad \eta + \xi - \xi\zeta = \eta + \xi - \eta\zeta = 0
\]
\[
\boxed{\zeta^2 = 0 \qquad \eta + \xi = \xi\zeta = \eta\zeta}
\]

\emph{NB}
\[
\det\sigma_{0V} = -\mu^2 = -1, \qquad
\det\sigma_{1V} = -(1+\zeta), \qquad
\det\sigma_{2V} = -1
\]
donc
\[
\boxed{(\det\sigma_{1V} = -1) \iff \zeta = 0} \ \Longrightarrow\ \xi = -\eta
\]

\page{66}
$\sigma_0$ est une \uncertain{symétrie par rapport à un hyperplan} sauf si $\mu = 1$ (\uncertain{cas d'un groupe de \ill{}})
\note{le début de la ligne est biffé ; la ligne est d'une écriture rapide et le complément entre parenthèses est à peine lisible.}

\[
\sigma_1(x,y,z) = (x,y,z) \iff -x + y + \xi z = 0,\ -x + y - \eta z = 0,\ \struck{(1+\zeta)}\,\zeta z = 0
\]
\uncertain{i.e. un hyperplan ssi} $\zeta = 0$, $\xi = -\eta$.

$s_1, s_2, s_3, s''_3$ coplanaires ssi somme des coord. de $s''_3$ $= 1$
\[
\Updownarrow
\]
\[
\xi + \eta \struck{+1+\zeta = 0} = \zeta \qquad (= \xi\zeta = \eta\zeta)
\]
il suffit $\zeta = 0$

On suppose par la suite $\boxed{\zeta = 0,\ \mu = 1}$ et on pose
\[
\begin{cases}
\lambda = 1 + \beta \\
\eta = 1 + \alpha
\end{cases}
\]

\[
\sigma_0 = \begin{pmatrix}
-1 & \beta & \beta & 1 \\
0 & 1 & 0 & 0 \\
0 & 0 & 1 & 0 \\
0 & 0 & 0 & 1
\end{pmatrix}
\qquad
\sigma_0(x,y,z) = \bigl(1 - x + \beta(y+z),\ y,\ z\bigr)
\]
\[
\sigma_1 = \begin{pmatrix}
0 & 1 & -(1+\alpha) & 0 \\
1 & 0 & 1+\alpha & 0 \\
0 & 0 & 1 & 0 \\
0 & 0 & 0 & 1
\end{pmatrix}
\qquad
\sigma_1(x,y,z) = \bigl(y - (1+\alpha)z,\ x + (1+\alpha)z,\ z\bigr)
\]
\[
\sigma_2 = \begin{pmatrix}
1 & 0 & 0 & 0 \\
0 & 0 & 1 & 0 \\
0 & 1 & 0 & 0 \\
0 & 0 & 0 & 1
\end{pmatrix}
\qquad
\sigma_2(x,y,z) = (x, z, y)
\]

\section{Formes quadratiques invariantes}

\page{67}
\[
\rho_s = \sigma_1\sigma_2 \qquad \rho_f = \sigma_0\sigma_1
\]
\note{ici $\rho_s$ et $\rho_f$ sont définis à l'inverse de la p.~\ref{page:61} ($\rho_s = \sigma_0\sigma_1$, $\rho_f = \sigma_1\sigma_2$) ; les indices sont repassés et c'est la lecture finale qui est donnée.}
\[
\rho_s^p = \mathrm{id} \iff F_p(\alpha) = 0 \qquad (\rho_s^d \neq \mathrm{id} \text{ si } 0 < d < p)
\]
\[
\rho_f^q = \mathrm{id} \iff F_q(\beta) = 0 \qquad (\rho_f^d \neq \mathrm{id} \text{ si } 0 < d < q)
\]
$F_p$ polynôme \uncertain{semi-cyclotomique}, $F_q$ id.

\emph{Formes quadratiques invariantes} \struck{en $x,y,z,t$}
\[
ax^2 + by^2 + cz^2 + uyz + vzx + wxy + px + qy + rz + d
\]
Conditions sur les coeff.\ ? Comme les \uncertain{éléments} \uncertain{fixent} $s_0$ \ill{}, \emph{OPS} $d = 0$.

L'invariance par $\sigma_0, \sigma_1$ dans le plan des $xy$ ($z = 0$) donne $a = b = -p = -q$, $w = -\beta a$, d'autre part l'inv.\ par $\sigma_2$ i.e.\ symétrie en $y, z$ donne $c = b$, $v = w$, $r = q$, de sorte que l'on doit avoir
\[
f(x,y,z) = a\bigl[x^2 + y^2 + z^2 \struck{\ill{}} - \beta x(y+z) - x - y - z\bigr] + uyz
\]
On va exprimer l'invariance par $\sigma_0$ puis $\sigma_1$, \ill{} qui \uncertain{va dire que} $f$ \uncertain{laisse} $uyz$ invariant \uncertain{par} $\sigma_0$, \ill{} il \uncertain{faudra} \struck{\ill{}} exprimer que

\struck{$a\bigl[x^2 + y^2 - \beta xy - x - y\bigr]$ est invariant, c'est le cas. Remarque : (\ill{} fonction \ill{} $z = 0$) \ill{}}
\struck{il faut exprimer l'invariance par $\sigma_1$, \ill{} invariante par $\sigma_0, \sigma_1$}

la condition \struck{$u =$}
\[
\struck{u = -a(\beta + 2(1+\alpha) + {}}
\]
\[
u = -a\bigl(2 + 2(\alpha+\beta) + \alpha\beta\bigr)
\]
d'où la forme de la \uncertain{fonction} \struck{quadratique} \uncertain{fondamentale}
\note{fin de page : la ligne « d'où la forme … » et le dernier mot sont d'une écriture très rapide.}

\page{68}
\[
f(x,y,z) = x^2 + y^2 + z^2 - \underbrace{\bigl(2 + 2(\alpha+\beta) + \alpha\beta\bigr)}_{\gamma = (\alpha+2)(\beta+2) - 2} yz - \beta zx - \beta xy - (x+y+z)
\]
Les fonctions \uncertain{semi}-quadratiques inv.\ par $\Gamma$ \uncertain{sont} donc les combinaisons linéaires de $f$ et de la constante 1. Celles qui s'annulent sur les sommets du polyèdre, i.e.\ en $s_0 = (0,0,0)$, sont les multiples de $f$.

\[
\begin{cases}
\text{Matrice de } f_0 \ (\text{quadratique}) = \begin{pmatrix} 2 & -\beta & -\beta \\ -\beta & 2 & -\gamma \\ -\beta & -\gamma & 2 \end{pmatrix} \\[4ex]
\delta'\ (\text{discrim.\ divisé}) = 4 - \beta^2(\gamma+2) - \gamma^2 = -(\alpha+2)(\beta+2)^2(\alpha-\beta) \\[1ex]
\gamma = 2(1+\alpha+\beta) + \alpha\beta
\end{cases}
\]
\marginal{Dans cas icos.\ $\alpha^2+\alpha-1 = 0$, $\beta = -1$, $\gamma = \alpha$, $\delta' = \emph{1}$ ok !}
\note{la parenthèse $(\alpha-\beta)$ est une lecture douteuse : le signe peut être $\alpha+\beta$.}
[inv.\ ssi $\alpha$ \uncertain{pas} $\neq \pm 2$, $\beta$ \uncertain{pas} $= -2$, $\alpha+\beta$ \uncertain{pas} $\neq 0$ ]
\marginal{2 cas.\ $\uparrow\beta, q$ \ill{}}

\[
\text{Matrice de } f = \begin{pmatrix}
2 & -\beta & -\beta & -1 \\
-\beta & 2 & -\gamma & -1 \\
-\beta & -\gamma & 2 & -1 \\
-1 & -1 & -1 & 0
\end{pmatrix}
\]
\[
\delta(f) = (\gamma+2)(\gamma - 4\beta - 6) = (\alpha+2)(\alpha-2)(\beta+2)^2
\]
il est \struck{\ill{}} ssi $\alpha$ \uncertain{pas} $\neq \pm 2$ i.e.\ \uncertain{car} $\neq p$, $\beta$ \uncertain{pas} $\neq -2$, $2\cdot\text{car} \neq q$
\uncertain{i.e.} $p$ = ordre des sommets, $q$ \ill{} faces \ill{}

\marginal{encerclé : $\delta(f) = -\alpha - 3$, \quad $N_{\mathbb{Q}(\alpha)/\mathbb{Q}}\,\delta(f) = 5$ !}
\marginal{note au crayon dans la marge gauche, montant le long du bord : \ill{}}

Linéaire et \uncertain{invariant} par $\sigma_0, \sigma_1, \sigma_2$, i.e.\ par $\Gamma$ \ill{}
\[
\sigma = \Bigl(\frac{\alpha}{2(\alpha+\beta)},\ \frac{-1}{2(\alpha+\beta)},\ \frac{-1}{2(\alpha+\beta)}\Bigr)
\]
\marginal{si $2(\alpha+\beta) \neq 0$ ; \uncertain{si $\alpha+\beta$ inv.}}
l'antipodisme est donné par
\[
\mathbf{a}(x,y,z) = \Bigl(\frac{\alpha}{\alpha+\beta} - x,\ \frac{-1}{\alpha+\beta} - y,\ \frac{-1}{\alpha+\beta} - z\Bigr)
\]
\note{la première coordonnée de $\sigma$ est écrite $\alpha$ par-dessus un premier $\alpha$ ; le « $\mathbf{a}$ » transcrit la lettre qu'il souligne.}

\page{69}
On vérifie que $f$ est invariant par $\mathbf{a}$, et \ill{} :
\[
\struck{f(\sigma) = \frac{-\alpha^2 - \alpha\beta + 2(\alpha+\beta)}{4(\alpha+\beta)^2}}
\qquad
\boxed{f(\sigma) = \frac{2-\alpha}{4(\alpha+\beta)}}
\]
C'est nul ssi $\alpha = 2$ (i.e., \add{en car.\ $p > 0$,} ssi la car.\ est égale à \struck{\ill{}} l'ordre \struck{\ill{}} des \struck{\ill{}} \ill{}, \uncertain{dans le cas réel}, la quadrique circonscrite est singulière\ldots)

\begin{itemize}
\item \emph{Cas tétraédral} \quad $\alpha = \beta = -1$, \quad $\delta' = 2$, \quad $\delta = -3$.

centre à l'$\infty$ ssi car.\ 2 ; $f_0$ singulière ssi car.\ 2 ; $f$ singulière ssi car.\ 3

$Q_E$ bisingulière ssi car.\ $\neq 2, 3$ (alors le centre est à dist.\ finie)
\item \emph{cube} \struck{\uncertain{cas octaédral}} \quad $\alpha = -1$, $\beta = 0$, \quad $\delta' = +4$, $\delta = -12$.

centre à l'$\infty$ ssi car.\ 2 ; $f_0$ singulière ssi car.\ 2 ; $f$ singulière ssi car.\ 2, 3

$Q_E$ \uncertain{bising.}\ ssi car.\ $\neq 2, 3$ (alors centre à dist.\ finie)
\item \emph{octaèdre} \struck{\ill{}} \quad $\alpha = 0$, $\beta = -1$, \quad $\delta' = 2$, $\delta = -4$.

centre à l'$\infty$ ssi car.\ 2 ; $f_0$ sing.\ ssi car.\ 2 ; $f$ sing.\ ssi car.\ 2

$Q_E$ \uncertain{bising.}\ ssi car.\ $\neq 2$ (alors centre à dist.\ finie)
\item \emph{icosaèdre} \quad $\alpha$ satisfait $\alpha^2 + \alpha - 1 = 0$, $\beta = -1$, \quad $\delta' = 1$ \struck{\ill{}}, $\delta = -\alpha - 3$, $N(\delta) = 5$.

centre à l'$\infty$ ssi car.\ 2 ; $f_0$ tjrs \uncertain{régulière} ; $f$ sing.\ ssi car.\ 5

$Q_E$ \uncertain{bising.}\ ssi car.\ $\neq 5$ (le centre \uncertain{jamais} à l'$\infty$ !)
\item \emph{dodécaèdre} \quad $\alpha = -1$, $\beta$ satisfait \struck{\ill{}} $\beta^2 + \beta - 1 = 0$,

\struck{$\delta' = \ldots\ N(\delta') = \ldots$ !}

$\delta' = \beta + 2$, $N(\delta') = 1$, \quad $\delta = -3(\beta+2)^2 = -9\beta - 15$, $N(\delta) = 9$.

centre à l'$\infty$ ssi car.\ 2 ; $f_0$ tjrs \uncertain{régulière} ; $f$ sing.\ ssi car.\ 3

$Q_E$ \uncertain{bising.}\ ssi car.\ $\neq 3$ (le centre \uncertain{jamais} à l'$\infty$)
\item \emph{pentagrammes} \quad $\alpha \neq \beta$ solutions de $\alpha^2 + \alpha - 1 = 0$, $\alpha\beta = -1$, $\beta = -1/\alpha = -1-\alpha$,

$\delta' = \beta + 2$, $N(\delta') = 1$, \quad $\delta = 4\alpha - 3$, $N(\delta) = 5$.

centre à l'$\infty$ ssi car.\ 2 ; $f_0$ tjrs \uncertain{régulière} ; $f$ bisingulière ssi car.\ 5

$Q_E$ \uncertain{bising.}\ ssi car.\ $\neq 5$ (le centre \uncertain{jamais} à l'$\infty$)
\end{itemize}
\note{dans le cas tétraédral, « car.\ $\neq 2,3$ » pour $Q_E$ est lu tel quel ; dans le cas pentagrammes, le $\delta'$ est écrit par-dessus une première valeur.}

\marginal{un gros « ? » en marge gauche}
Il y a \struck{\ill{}} un \struck{\ill{}} \uncertain{cas} \add{fini} où r.l.\ \uncertain{rég.}\ \uncertain{est} \emph{bisingulier}, mais \uncertain{je ne} \ill{} pas \ill{} \ill{} en car.\ 3 \struck{\ill{}} (cube), 3 (icos.), 5 (dodéc.).

\page{70}
\[
\boxed{F(X,Y,Z) = (-X+Y+Z)^2 - 4(\alpha+2)YZ + \underbrace{\frac{2-\alpha}{4(\alpha+2)}}_{\lambda > 0}}
\]
\struck{(\uncertain{origine} \ill{})}
\[
-2 < \alpha < 2, \qquad 0 < 4(\alpha+2) < 16
\]
\[
4YZ = (Y+Z)^2 - (Y-Z)^2
\qquad
\begin{cases}
-X + Y + Z = U\sqrt{\lambda} \\
\sqrt{\alpha+2}\,(Y+Z) = W\sqrt{\lambda} \\
\sqrt{\alpha+2}\,(Y-Z) = V\sqrt{\lambda}
\end{cases}
\]
\note{les lettres de droite sont surchargées (les $U$, $V$, $W$ ont été réattribués) ; le système est donné dans son état final probable.}
\[
F_1(U,V,W) = U^2 \struck{+ \ldots}\ (V^2 - W^2) + \ldots\ 1
\]
\note{une ligne largement raturée, que l'on ne peut rétablir en entier.}

\[
U^2 + V^2 - W^2 + \struck{\lambda} = 0, \qquad U^2 = V^2 + W^2 + \struck{\lambda}
\]
\[
\struck{W^2 - \rho UV + \sigma = 0} \qquad \struck{\rho UV = W^2 + \sigma}
\]
\[
\struck{W^2 - \rho YZ + \sigma = 0} \qquad \struck{\rho YZ = W^2 + \sigma}
\]
\[
W^2 = U^2 + V^2 + 1
\]
\drawing{Deux croquis d'axes : à gauche, un trièdre $U, V, W$ avec un plan hachuré et des traces de la quadrique ; à droite, un trièdre $u, V, W$ seul. Des courbes (hyperboles, une sinusoïde tracée à travers les formules biffées) accompagnent.}

\page{71}
\note{la moitié supérieure de la page est barrée de grandes croix.}
\struck{Nous allons examiner de plus près le cas où $\alpha, \beta$ sont des racines de polynômes \uncertain{semi-}cyclotomiques $F_p(\alpha) = 0$, $F_q(\beta) = 0$ ($p, q \geq 1$, sans exclure les cas 1, 2 qui correspondent \add{\uncertain{resp.}} aux valeurs $2, -2$ pour $\alpha, \beta$\ldots), \ill{} le cas \emph{réel} $k = \mathbb{R}$. Si on veut obtenir une \struck{\ill{}} \ill{} de la réalisation topologique de la cette combinatoire $C_{p,q}$, il faut prendre (pour $p$ resp.\ $q$ fini $\geq 3$) des valeurs de $\alpha$ resp.\ $\beta$ de la forme $\zeta + \zeta^{-1}$, avec $\zeta = \exp 2i\pi/n$ ($n = p$ resp.\ $q$) ; d'autre part, $\alpha = -2$ \add{ou $\beta = -2$} donne des croisements, donc si $p = \infty$ il faut $\alpha = 2$, et $q = \infty$ \uncertain{il} faut $\beta = 2$.}

\[
\beta = 2 \qquad \alpha \neq -2
\]
\[
\gamma = 2(3+\alpha) + 2\alpha = 2(2\alpha+3) \qquad \gamma + 2 = 2(2\alpha+4) = 4(\alpha+2)
\]
\[
\delta' = 4 - 4(\gamma+2) - \gamma^2 = -(\gamma^2 + 4\gamma + 4) = -(\gamma+2)^2 = -16(\alpha+2)^2
\]
qui \uncertain{s'annule} que si $\alpha = -2$, cas que nous avons exclu\ldots
\[
f_0(x,y,z) = x^2 + y^2 + z^2 - 2(2\alpha+3)yz - 2zx - 2xy = (-x+y+z)^2 - 2(2\alpha+4)yz
\]
\[
f_0(x,y,z) = (-x+y+z)^2 - 4(\alpha+2)yz
\]
\note{les formules en $x^2+y^2+z^2 - 2(2\alpha+3)yz - 2zx - 2xy$ supposent $\beta = 2$ dans la forme de la p.~\ref{page:68} ; le coefficient de $yz$ y est écrit $-2(2\alpha+3)$, alors que $\gamma = 2(2\alpha+3)$ : tel quel.}

\section{Calculs en coordonnées universelles (icosaèdre)}
\note{titre pris de la chemise p.~73, de sa main : « (32) Calculs en coordonnées universelles (icosaèdre) — vieilles rédactions ».}

\page{74}
\drawing{Icosaèdre en perspective : au centre le triangle $s_0 s_1 s_2$ (avec $\sigma_1$, $\sigma_0$, $\sigma_2$ marqués sur ses droites), autour les sommets $s_3$, $s'_3$, $s''_3$ et les antipodes $\mathbf{a}s_0$, $\mathbf{a}s_1$, $\mathbf{a}s_2$, $\mathbf{a}s_3$, $\mathbf{a}s'_3$, $\mathbf{a}s''_3$ ; arêtes cachées en pointillé.}

\[
\sigma_0 : \begin{cases} s_0 \leftrightarrow s_1 \\ s_2 \mapsto s_2 \\ s_3 \mapsto s_3 \end{cases}
\qquad
\begin{cases} e_1 \mapsto -e_1 \\ e_2 \mapsto -e_1 + e_2 \\ e_3 \mapsto -e_1 + e_3 \end{cases}
\qquad
\sigma_0 = \begin{pmatrix} -1 & -1 & -1 & 1 \\ 0 & +1 & 0 & 0 \\ 0 & 0 & +1 & 0 \\ 0 & 0 & 0 & 1 \end{pmatrix}
\]
\[
\boxed{\sigma_0(x,y,z) = \bigl(1 - (x+y+z),\ y,\ z\bigr)}
\]

\[
\sigma_1 : \begin{cases} s_0 \mapsto s_0 \\ s_1 \mapsto s_2 \\ s_2 \mapsto s_1 \\ s_3 \mapsto s'_3 = s_0 + (1+\alpha)(e_2 - e_1) + e_3 \end{cases}
\]
\[
\begin{cases} e_1 \mapsto e_2 \\ e_2 \mapsto e_1 \\ e_3 \mapsto (1+\alpha)(e_2 - e_1) + e_3 = \alpha' e_1 - \alpha' e_2 + e_3 \end{cases}
\]
\[
\sigma_1 = \begin{pmatrix} 0 & 1 & \alpha' & 0 \\ 1 & 0 & -\alpha' & 0 \\ 0 & 0 & 1 & 0 \\ 0 & 0 & 0 & 1 \end{pmatrix}
\qquad
\boxed{\sigma_1(x,y,z) = (y + \alpha' z,\ x - \alpha' z,\ z)}
\]
\note{$(1+\alpha)(e_2 - e_1) = \alpha' e_1 - \alpha' e_2$ suppose $\alpha' = -(1+\alpha)$, ce que la p.~\ref{page:80} écrit $\alpha + \alpha' = -1$.}

\drawing{Seconde figure : l'icosaèdre projeté dans le triangle $\bar s_1, \bar s_0, \bar s_2$, avec au centre $s_0, s_1, s_2, s_3, s'_3, s''_3$, $\bar s_3$, $\bar s'_3$, $\bar s''_3$ ; esquisse au crayon dessous.}

\[
\sigma_2 : \begin{cases} s_0 \mapsto s_0 \\ s_1 \mapsto s_1 \\ s_2 \mapsto s_3 \\ s_3 \mapsto s_2 \end{cases}
\qquad
\begin{cases} e_1 \mapsto e_1 \\ e_2 \mapsto e_3 \\ e_3 \mapsto e_2 \end{cases}
\qquad
\sigma_2 = \begin{pmatrix} 1 & 0 & 0 & 0 \\ 0 & 0 & 1 & 0 \\ 0 & 1 & 0 & 0 \\ 0 & 0 & 0 & 1 \end{pmatrix}
\qquad
\boxed{\sigma_2(x,y,z) = (x, z, y)}
\]

\[
\begin{cases}
H_0 : 2x + y + z - 1 = 0 \\
H_1 : x - y - \alpha' z = 0 \\
H_2 : y - z = 0
\end{cases}
\qquad
\begin{cases}
\Delta_0 : y = z = 0 & \text{i.e. } \Delta_0 = \mathcal{O}_S\cdot e_1 \\
\Delta_1 : z = 0,\ x + y = 0 & \text{i.e. } \Delta_1 = \mathcal{O}_S(e_1 - e_2) \\
\Delta_2 : x = 0,\ y + z = 0 & \text{i.e. } \Delta_2 = \mathcal{O}_S(e_2 - e_3)
\end{cases}
\]
en car.\ 2, $H_2 \parallel H_0$.
\note{l'équation de $H_0$ porte une surcharge devant le « $1$ » ; $H_2 \parallel H_0$ en car.\ 2 ne se lit pas sur les équations transcrites, qui sont peut-être inexactes à cet endroit.}

\[
\begin{cases}
D_0 = H_1 \cap H_2 : y = z = \alpha' x \\ \qquad D_0 = \mathcal{O}_S(1, \alpha', \alpha') = \mathcal{O}_S(-\alpha, 1, 1) \\
D_1 = H_2 \cap H_0 : y = z = \tfrac12 - x \ (\text{car.} \neq 2) \\ \qquad \mathrm{hom}\,D_1 : (y = z = -x) = \mathcal{O}_S(1, -1, -1) \\
D_2 = H_0 \cap H_1 : y = (1-3\alpha)x + \alpha,\ z = \ldots(1-\alpha)x + (1-\alpha) \\ \qquad \mathrm{hom}(D_2) = \mathcal{O}_S\bigl(1, 1-3\alpha, \ldots(1-\alpha)\bigr)
\end{cases}
\]
\note{les coefficients de $D_2$ sont surchargés ; ce qui est donné est la lecture la plus probable, avec des facteurs initiaux illisibles marqués « $\ldots$ ».}

\emph{centre} \quad $c = \tfrac12(\alpha', 1-\alpha', 1-\alpha') = \struck{\ill{}}$

\emph{Antipodisme} \quad $\mathbf{a}(x,y,z) = \bigl(\alpha' - x,\ (1-\alpha') - y,\ (1-\alpha') - z\bigr)$
\note{sous cette ligne, deux restes de calcul : « $\tfrac12(-1+\sqrt5)$ » et « $\lambda = -2 + \sqrt8$ » ; le premier $\tfrac12$ est lu \uncertain{$\tfrac12$}.}

\page{75}
\[
\begin{array}{ll}
s_0 = (0,0,0) & \mathbf{a}s_0 = (\alpha', 1-\alpha', 1-\alpha') \\
s_1 = (1,0,0) & \mathbf{a}s_1 = (\alpha'-1, 1-\alpha', 1-\alpha') \\
s_2 = (0,1,0) & \mathbf{a}s_2 = (\alpha', -\alpha', 1-\alpha') \\
s_3 = (0,0,1) & \mathbf{a}s_3 = (\alpha', 1-\alpha', -\alpha') \\
s'_3 = (\alpha', -\alpha', 1) & \mathbf{a}s'_3 = (0, 1, -\alpha') \\
s''_3 = (0, -\alpha', 1) & \mathbf{a}s''_3 = (\alpha', 1-2\alpha', -\alpha')
\end{array}
\]
\[
\begin{array}{l}
\mathbf{a}(x,y,z) = \bigl(\alpha' - x,\ (1-\alpha') - y,\ (1-\alpha') - z\bigr) \\[1ex]
c = \tfrac12(\alpha', 1-\alpha', 1-\alpha') \quad (\text{centre}) \\[1ex]
u = 2c = (\alpha', 1-\alpha', 1-\alpha', 2)
\end{array}
\]
\note{dans $\mathbf{a}s'_3$, la deuxième coordonnée est écrite « $1$ » après une rature ; le calcul direct donnerait $1-\alpha'+\alpha'$.}

\[
q_0(x,y,z) = x^2 + y^2 + z^2 - \alpha yz + x(y+z) - (x+y+z)
\]
\[
\Bigl[\ = (\text{en car.} \neq 2)\ \alpha\bigl[(x - \alpha y)^2 + \alpha'(x-\alpha y)(x - \alpha z) + (x - \alpha z)^2\bigr] - (x+y+z)\ \Bigr]
\]
Si $e$ est un vecteur arête (\uncertain{tel que} $s_1 - s_0 = e_1$), $f$ \uncertain{un vecteur} coarête (tel que $s_3 - s_2 = e_3 - e_2$)
\[
\begin{cases}
q(e) = 1 \quad (\text{c'est la normalisation de } q\,!) \\
q(f) = \alpha + 2 \\
q(f)/q(e) = \alpha + 2 = (\alpha+1)^2 = \alpha'^2
\end{cases}
\]

En coordonnées distinguées, si $(\pm x, \pm y, 0)$, $(\pm y, 0, \pm x)$, $(0, \pm x, \pm y)$ sont les 12 sommets de l'icosaèdre, on aura (si $(x,y,0)$ et $(-x,y,0)$ sont les sommets d'une arête, donc $(x,y,0)$ et $(x,-y,0)$ ceux d'une coarête)
\[
\struck{\ldots} = \frac{q(f)}{q(e)} = \frac{4y^2}{4x^2} = \Bigl(\frac{y}{x}\Bigr)^2 = \alpha'^2 \quad \text{donc} \quad \frac{y}{x} = \pm\alpha'
\]
(dans le cas réel, \uncertain{icosaèdre} convexe, $\alpha = \tfrac12(-1+\sqrt5)$, $\alpha' = \tfrac12(-1-\sqrt5)$ et si on prend $x, y > 0$, on trouve $y/x = -\alpha' = \tfrac12(1+\sqrt5)$)
\drawing{Petit croquis en marge : le rectangle d'or de sommets $s_1$, $s_0$, $-s_0$, $-s_1$, côtés $x$ et $y$.}

\[
\delta'(q_0) = 1 \qquad \text{Matrice } \varphi_{q_0} = \begin{pmatrix} 2 & 1 & 1 \\ 1 & 2 & -\alpha \\ 1 & -\alpha & 2 \end{pmatrix}
\]
Matrice non singulière sauf en car.\ 2, où le noyau de $\varphi_{q_0}$ est engendré par le vecteur $u = (\alpha', 1-\alpha', 1-\alpha') = (\alpha', \alpha, \alpha) = (1+\alpha, \alpha, \alpha)$
(\emph{NB} $q_0(u) = \alpha' \neq 0$, ce qui exprime la non-singularité de $q_0$ en car.\ 2)
\note{« $(\alpha', \alpha, \alpha) = (1+\alpha, \alpha, \alpha)$ » vaut en car.\ 2 seulement, où $\alpha' = 1+\alpha$.}

\page{76}
\emph{Vecteurs remarquables} \add{et directions remarquables} (\uncertain{cas} car.\ $\neq 2$)

On suppose l'icosaèdre \uncertain{à centre $= 0$}, identifiant $E$ à l'espace des translations $V$. \uncertain{Utilisant} la forme quadratique, on identifie aussi $V$ et $\check V$ --- donc $\check{\mathbb{P}}(V)$ (directions de droites dans $E$) et $\mathbb{P}(V)$ (directions de plans dans $E$).

\begin{itemize}
\item $S$ --- 12 vecteurs « \add{doubles} sommets » \uncertain{$s_i - \mathbf{a}s_i$}
\item $\Delta$ --- 30 vecteurs \add{arêtes} \struck{\ill{}} $s_i - s_j$ ($\{i,j\} \in A$) se groupent en 15 paires de vecteurs opp.
\marginal{\emph{NB} chaque arête \ill{} un vecteur, mais deux arêtes \ill{} \uncertain{doivent} définir le même \ill{} \struck{les vecteurs}}
\item $\Delta'$ --- 30 vecteurs coarêtes $s_i - s_j$ ($\{i,j\} \in A'$), ce sont les \ill{} ($a \in \Delta'$)
\[
\struck{\ldots} \qquad \boxed{\Delta' = \alpha'\Delta, \quad \Delta = \alpha\Delta'}
\]
\drawing{petit rectangle à flèches, figurant une arête et une coarête.}
\item $\mathcal{A}$ --- Les \uncertain{positions} \add{30} d'arêtes $\tfrac12(s_i + s_j) - 2c$ ($\{i,j\} \in A$)
\item $\mathcal{A}'$ --- \add{30} de coarêtes $\tfrac12(s_i + s_j)$ ($\{i,j\} \in A'$)
\[
\Delta = \struck{\ldots}\,\mathcal{A}', \quad \Delta' = \struck{\ldots}\,\mathcal{A}
\qquad
\begin{cases} \mathcal{A} = \alpha'\mathcal{A}' \\ \mathcal{A}' = \alpha\mathcal{A} \end{cases}
\]
i.e.\ $\Delta$ est aussi formé des $s_i + s_j$ ($\{i,j\} \in A'$), $\Delta'$ \uncertain{des mêmes} ($\{i,j\} \in A$)
\item $\Phi$ --- Les \add{20, triples} barycentres des faces $\uncertain{\tfrac13}(s_i + s_j + s_k) - \ldots c$ ($\{i,j,k\} \in F$)
\item $\Phi'$ --- Les \add{20, triples} barycentres des \uncertain{cofaces} ($\{i,j,k\} \in F'$)
\[
\Phi' = \lambda\Phi, \quad \Phi = \lambda'\Phi' \qquad (\lambda\lambda' = \mathrm{Norm}(\lambda) = 1) \qquad \lambda = 2\alpha - 1
\]
($\Phi$ et $\Phi'$ sont respectivement les sommets des dodécaèdres inscrits dans $\Pi$ et dans le \uncertain{cofacettoïde associé} $\Pi'$)
\item $\Psi$ --- Les 20 sommets \add{multiplicité un 10} des dodécaèdres circonscrits \add{à $\Pi$}, \uncertain{intersections} des plans des faces adjacentes : \ill{}
\item $\Psi'$ --- item pour $\Pi'$
\[
\boxed{\Psi' = \mu\Psi, \ \ \Psi = \mu'\Psi'} \qquad \struck{\mu\mu' = \mathrm{Norm}(\mu) = 1} \quad \gamma = -\lambda
\]
\[
\boxed{\Psi = \mu\Phi, \ \ \Psi' = \mu'\Phi'} \qquad \mu = 1 + 2\alpha = -(1+2\alpha') \quad (\mu\mu' = -5),\ \mu^2 = 5 \text{ i.e. } \mu = \sqrt5
\]
\item $\Sigma$ --- Les 12 \add{\uncertain{multiples}} barycentres des pentagones \uncertain{inscrits} aux sommets $2(s_1 + s_2 + s_3 + s_4 + s_5) - 2c$
\[
\Sigma = \rho S = \rho' S \qquad \rho' = -\rho \quad \text{i.e. } \mathrm{Tr}\,\rho = 0 \quad \struck{\ldots} \qquad \rho = \mu = 1 + 2\alpha\ldots
\]
\end{itemize}
\note{page d'une écriture très rapide : les listes de vecteurs sont sûres dans leur structure (lettres de gauche, cardinaux, relations $\Delta' = \alpha'\Delta$, $\Phi' = \lambda\Phi$, $\Psi = \mu\Phi$, $\Sigma = \rho S$), beaucoup moins dans leurs légendes.}

\page{78}
Directions remarquables
\[
\begin{array}{rl}
6 & \text{directions de sommets} \\
15 & \text{directions d'arêtes} \\
10 & \text{directions de plans} \\ \hline
31 & = 5^2 + 5 + 1 \text{ nb de pts du plan projectif sur } \mathbb{F}_5
\end{array}
\]
Les 6 directions de sommets correspondent à la conique, les \add{15} directions d'\add{arêtes} \struck{plans} correspondent aux sécantes \uncertain{divisant un diviseur} \add{\uncertain{discriminant}} $/\mathbb{F}_5$, les \add{10} directions de plans correspondent aux diviseurs \uncertain{conjugués} $/\mathbb{F}_5$.

vecteurs remarquables en \emph{car.\ 5}
\[
1 + 12 + 60 + 40 = 113 \qquad 125 = 1 + \underbrace{24}_{\substack{\text{pts sur la conique}\\ = 2\times 12\\ (\text{deux icos.})}} + 60 + 40
\]
\struck{les \add{25} directions \uncertain{non isotropes} \ill{} associées aux \ill{} \ill{} des \uncertain{tridièdres} dist.\ \ill{}}

\struck{20 aux \ill{} \ill{} dans \ill{}}
\note{« $113$ » est surchargé (peut-être écrit par-dessus un premier total).}

\drawing{Croquis : un parallélogramme de sommets $s_0$, $s'_1$, $s_1$, $s_0$ autour d'un point $c$, deux vecteurs issus de $c$ ; à côté, « $x^2+y^2$ » biffé.}

\page{80}
\note{la page est coupée sur sa droite dans la numérisation : les fins des lignes de la colonne de droite manquent.}
\[
y = z \qquad x = (1+\alpha')y = -\alpha y \qquad y = \alpha' x
\]
\[
(-\alpha y, y, y) \text{ ou } (x, \alpha' x, \alpha' x), \qquad (-\alpha, 1, 1) \qquad (1, \alpha', \alpha')
\]
\[
-\alpha = \alpha'(1-\alpha') = \alpha' - \alpha'^2 \qquad \alpha' = -\alpha(1-\alpha') = -\alpha - 1 \ \text{ok}
\]
\[
\mathcal{O}_S(-\alpha, 1, 1)
\qquad
\begin{array}{l} y + z + 1 = 0 \\ x + y + \alpha z = 0 \\ y = z \end{array}
\]
$e_1 \mapsto \ldots$
\[
\begin{cases}
\bar s_0 = (-\alpha, 1, 1) \\
\bar s_1 = \sigma_0 \bar s_0 = (\alpha - 1, 1, 1) \\
\text{d'où } \bar s_1 - \bar s_0 = (2\alpha - 1, 0, 0), \quad \alpha' = 1 - \alpha'
\end{cases}
\qquad \alpha + \alpha' = -1
\]
\note{colonne de droite, tronquée : « $1 - (2 - \ldots$ », « $\alpha - 1$ », « $\alpha = \tfrac12 \ldots$ », « $\tfrac14 + \ldots$ », « $2\alpha - 1 = 2\alpha + \ldots$ ».}

\[
f(0) = f(\sigma_1\sigma_2) = 0
\qquad
f(\sigma_0^2) = f(\sigma_0) + \sigma_0 f(\sigma_0) \ldots
\qquad e + \sigma_0 e = 0, \quad \boxed{\sigma_0 e = -e}
\]
\[
\begin{cases}
\Gamma \to SO(q) \\
+ \text{ 1-cocycle } \Gamma \xrightarrow{\ f\ } E
\end{cases}
\qquad
f(gg') = f(g) + g f(g')
\]
\uncertain{Il suffit} de \uncertain{déterminer}
\[
f(\sigma_1) = f(\sigma_2) = 0 \qquad f(\sigma_0) = \struck{\ldots}
\]
\[
\sigma_0\sigma_2\sigma_0\sigma_2 \qquad f(\sigma_0\sigma_2) = \struck{\ldots}\, e \qquad e + (\sigma_0\sigma_2)e = \ldots \qquad \struck{\sigma_0\sigma_2 e = \ldots} \qquad \sigma_0(\sigma_2 e - e) \ldots
\]
\[
\sigma_0\sigma_1 = u \qquad f(u) = e,\ f(u^2) = e + ue,\ \ldots
\qquad
\frac{e_1 \wedge e_2 \wedge e_3}{q}
\]
\[
f(u^3) = f(u \cdot u^2) = e + u(e + ue) = e + ue + u^2 e
\]
\[
s_0 \mapsto s_1, \qquad s_2 - s_0 = e_2 - e_3 = (0, 1, -1)
\qquad
\boxed{q_0(e) = 1}
\]
\[
1 + 1 + \alpha = 2 + \alpha = \struck{\ldots}\ 1 - \alpha' = \boxed{\alpha'^2}
\]
\[
\boxed{\begin{cases} \sigma_0 e = -e \\ \sigma_2 e = e \end{cases}}
\iff (1 + \sigma_0)e = 0
\iff (1 + \sigma_0\sigma_2)e = 0
\qquad
\boxed{(1 + u + u^2)e = 0}\ \struck{\ldots}
\]
\note{« $s_2 - s_0 = e_2 - e_3$ » est tel quel sur la page (le $s_2$ est peut-être $s_3 - s_2$). Le calcul du 1-cocycle se poursuit vraisemblablement au-delà de ce lot.}

\end{document}
